% Generated draft: requires mathematical and source review.
\documentclass[11pt]{article}
\usepackage[margin=25mm]{geometry}
\usepackage{fontspec}
\setmainfont{Latin Modern Roman}
\usepackage{amsmath,amssymb,mathtools,unicode-math}
\setmathfont{Latin Modern Math}
\usepackage{enumitem,needspace,xurl,hyperref}
\setlength{\parindent}{0pt}
\setlength{\parskip}{5pt}
\setlength{\emergencystretch}{3em}
\newcommand{\N}{\mathbb{N}}
\newcommand{\Z}{\mathbb{Z}}
\newcommand{\Q}{\mathbb{Q}}
\newcommand{\R}{\mathbb{R}}
\newcommand{\C}{\mathbb{C}}
\newcommand{\F}{\mathbb{F}}
\newcommand{\A}{\mathbb{A}}
\newcommand{\PP}{\mathbb{P}}
\newcommand{\E}{\mathbb{E}}
\newcommand{\eps}{\varepsilon}
\newcommand{\dd}{\,\mathrm{d}}
\begin{document}
\section*{UnsolvedMath ID: 700023}
\textbf{Source problem code:} catalogue entry
\subsection{Definitions and mathematical statement}
Work in the first-order language of groups \(\mathcal L_{\mathrm{grp}}=\{1,\cdot,{}^{-1}\}\), without parameters. A universal sentence is a sentence \(\forall x_1\cdots\forall x_n\,\varphi(x_1,\ldots,x_n)\), where \(n\geq 0\) is finite and \(\varphi\) is a quantifier-free Boolean combination of equations between group terms. For a group \(G\), define its universal theory by
\[
\operatorname{Th}_{\forall}(G)=\{\sigma:\sigma\text{ is a universal }\mathcal L_{\mathrm{grp}}\text{-sentence and }G\models\sigma\}.
\]
For a set \(X\), let \(F_X\) be the free group on \(X\). Define \(\gamma_1(F_X)=F_X\) and \(\gamma_{j+1}(F_X)=[\gamma_j(F_X),F_X]\), where \([H,K]\) denotes the subgroup generated by \([h,k]=h^{-1}k^{-1}hk\) with \(h\in H\) and \(k\in K\). For a positive integer \(c\), the free nilpotent group on \(X\) in the variety of groups of nilpotency class at most \(c\) is
\[
N_{X,c}=F_X/\gamma_{c+1}(F_X).
\]
\textbf{Problem (M. I. Kargapolov).} Describe the universal theory of a free nilpotent group; in the notation above, determine \(\operatorname{Th}_{\forall}(N_{X,c})\).

\textit{Unresolved scope.} The supplied statement does not specify the rank \(|X|\), whether infinite ranks or the trivial group are included, or the nilpotency parameter \(c\). The displayed construction makes these parameters explicit without selecting their intended range. The source also does not specify whether ``describe'' requires an explicit axiomatization, an effective membership criterion, or another characterization, nor whether the description must be uniform in rank and class.
\subsection{Short English statement}
Describe exactly which universally quantified, parameter-free statements in the language of groups hold in a free nilpotent group. The supplied problem leaves the rank, nilpotency class, and required form of the description unspecified.
\subsection{Sources}
\begin{itemize}
\item Local UnsolvedMath registry, permanent ID \(700023\), source code \texttt{catalogue entry}, title \emph{Kourovka Notebook Problem 1.28}; supplied upstream revision \texttt{d371ee917ac9647fbb06cd6fc529deec6f3dda82}. The record labels the problem open; this is supplied registry status, not an independently verified mathematical or literature assessment.
\item \emph{The Kourovka Notebook (Unsolved Problems in Group Theory)}, original Problem 1.28, first issue (1965); supplied edition reference: 21st issue (2026), \texttt{arXiv:1401.0300v48}. The supplied collection description names E. I. Khukhro and V. D. Mazurov as editors. This reference is supplied attribution and has not been independently consulted.
\item Proposed by M. I. Kargapolov, with proposed year 1965 according to the supplied record.
\end{itemize}
\end{document}
